% ===== latex-recomposed:document =====
% flm: style file flmrecomposed.5bb7fbac.sty
\flmHeading[toc=true]
\section{Statements}{\flmLabels{}}

\begin{theorem}[Main result]\flmLabels{\label{ref1}}
Every shipped kind is written in this document.
\end{theorem}

\begin{proof}[{\flmProofTitle{\flmRef{ref1}}}]
By \flmRef{ref2} and \flmRef{ref3}; together, \flmRef{ref1,ref2}.
\end{proof}

\begin{lemma}[Key lemma]\flmLabels{\label{ref2}}
A lemma proved below, its proof naming it.
\end{lemma}

\begin{proof}[{\flmProofTitle{\flmRef{ref2}}}]
One star shows the statement proved.
\end{proof}

\begin{lemma}\flmLabels{\label{ref3}}
A lemma whose proof records the relation without showing it.
\end{lemma}

\begin{proof}
Two stars record the relation only.
\end{proof}

\begin{proposition}\flmLabels{\label{ref4}}\noproofref A proposition whose proof is not coming.
\end{proposition}

\begin{corollary}\flmLabels{\label{ref5}}
\end{corollary}

\begin{proof}[{\flmProofTitle{\flmRef{ref5}}}]
The corollary above has an empty body.
\end{proof}

\begin{conjecture}\flmLabels{\label{ref6}}
A conjecture expects no proof.
\end{conjecture}

\flmHeading[toc=true]
\section{Definitions and the other kinds}{\flmLabels{}}

\begin{definition}[Opening with a list]\flmLabels{\label{ref7}}
\begin{itemize}
\item a body opening with a list gives the heading a paragraph of its own;
\item and the list follows it.
\end{itemize}
\end{definition}

\begin{remark}\flmLabels{\label{ref8}}
A remark.
\end{remark}

\begin{idea}
An idea.
\end{idea}

\begin{question}
A question?
\end{question}

\begin{claim}
A claim.
\end{claim}

\begin{observation}
An observation.
\end{observation}

\begin{problem}
A problem.
\end{problem}

\begin{exercise}[Host's kind]\flmLabels{\label{ref9}}
An exercise, a kind the host configures.
\end{exercise}

See \flmRef{ref9}, \flmRef{ref7} and \flmRef{ref8}.

\begin{itemize}
\item\begin{lemma}\flmLabels{\label{ref10}}
A lemma in a list item.
\end{lemma}
\item\begin{proof}[{\flmProofTitle{\flmRef{ref10}}}]
Its proof in the next item.
\end{proof}
\end{itemize}

A note holds a remark.\flmXNfootnote{\begin{remark}\flmLabels{\label{ref11}}
A remark in a footnote.
\end{remark}} See \flmRef{ref11} and \flmRef{ref10}.

\begin{flmUEkeptthm}{Kept}
A theorem inside a user definition's template.
\end{flmUEkeptthm}%%
% ===== latex-recomposed:style flmrecomposed.5bb7fbac.sty =====
\NeedsTeXFormat{LaTeX2e}
\ProvidesPackage{flmrecomposed.5bb7fbac}
% --- packages ---
% flm.latex.pkg.amsthm
\RequirePackage{amsthm}
% flm.latex.pkg.zref-clever
\RequirePackage{zref-clever}
% --- commands and environments the library defines for the document ---
% flm.latex.def.xnotes.footnote
% \flmXNfootnote[<labels>]{<note>}: a note of the kind `footnote' in the FLM document;
% by default it passes the note, its labels first, to \footnote.
% Redefine it to set these notes another way.
\ProvideDocumentCommand\flmXNfootnote{+O{} +m}{\footnote{#1#2}}
% flm.latex.def.theorems
% Theorems. \noproofref, written after a statement's labels, says that no proof
% of the statement is coming; defined here as nothing where no package (phfthm)
% defined it first.
\ProvideDocumentCommand\noproofref{}{}
% The kinds of a class are set in the theorem style of the class's name. The
% shipped class definitionlike is set in amsthm's definition style (upright);
% declare a theorem style of that name before loading this file to change it.
\ifcsname th@definitionlike\endcsname\else\newtheoremstyle{definitionlike}{}{}{\normalfont}{}{\bfseries}{.}{5pt plus 1pt minus 1pt}{}\fi
% flm.latex.def.theorems.resolved
% \flmProofTitle{<title>}: what a proof's heading sets for a title -- one the
% document wrote, or the statement the proof refers back to -- which is written
% as its argument inside the proof environment's optional argument, the whole
% heading amsthm sets. By default Proof (<title>), \proofname being the word
% amsthm defines. Redefine \flmProofTitle to set a proof's title differently;
% \renewcommand\flmProofTitle[1]{#1} sets the title alone.
\ProvideDocumentCommand\flmProofTitle{m}{\proofname\space (#1)}
% flm.latex.def.refs
% \flmRef[<options>]{<label>}: a reference that gives the name of its target
% ("Section 2"), set by \zcref of the zref-clever package.
%
% The options are the ones the document wrote, in FLM's own vocabulary. This
% definition translates each of them by its key -- a value written after the
% key is ignored, as in FLM -- and drops every key it has no translation for:
%
%   noprefix    -> noname                 the number or title alone, with no
%                                         name before it ("2", not
%                                         "Section 2")
%   S           -> S                      the first name capitalized and
%                                         unabbreviated, for a reference that
%                                         opens a sentence
%   capital     -> cap=true,abbrev=false  the names capitalized ("Lemma 1")
%   lowercase   -> cap=false,abbrev=false the names in lowercase ("lemma 1")
%   abbreviated -> abbrev=true            the abbreviated names ("Lem. 1")
%                                         where the reference type declares
%                                         them, the full names otherwise
%   anything else -> dropped              the default wording; a variant of
%                                         your own naming has no counterpart
%                                         in zref-clever, which would stop on
%                                         a key it does not know
%
% Redefine \flmRef to set references differently -- another package, another
% wording, your own variants, no link. Your definition sees the options as the
% document wrote them, in FLM's vocabulary, and maps them as it likes.
%
% The translation uses expl3, LaTeX's programming layer, which is part of the
% format and needs no package; its names read only between \ExplSyntaxOn and
% \ExplSyntaxOff. \keyval_parse:NNn reads the option list as LaTeX reads any
% key-value list: it hands a key written alone to the first macro below and a
% key with its value to the second, which drops the value. \str_case:nnF looks
% the key up in the table and appends what \zcref is to be given, or nothing.
% The macros with @ in their names are internals, not a customization point:
% redefine \flmRef, not its parts.
\ExplSyntaxOn
\def\flm@ref@add#1{\tl_put_right:Nn \flm@ref@options {#1,}}
\def\flm@ref@key#1{\str_case:nnF {#1}
  {{noprefix}{\flm@ref@add{noname}}
   {S}{\flm@ref@add{S}}
   {capital}{\flm@ref@add{cap=true,abbrev=false}}
   {lowercase}{\flm@ref@add{cap=false,abbrev=false}}
   {abbreviated}{\flm@ref@add{abbrev=true}}}
  {}}
\def\flm@ref@pair#1#2{\flm@ref@key{#1}}
\def\flm@ref@zcref#1#2{\zcref[#1]{#2}}
\ProvideDocumentCommand\flmRef{O{} m}{%
  \def\flm@ref@options{}%
  \keyval_parse:NNn \flm@ref@key \flm@ref@pair {#1}%
  \exp_args:No \flm@ref@zcref {\flm@ref@options} {#2}}
\ExplSyntaxOff
% flm.latex.def.label-display-text
% \flmLabelsWithDisplayText{<text>}{<labels>}: makes <labels> refer to this
% place, with <text> as what a reference to them prints. It is the only
% command for that: a heading the class did not number uses it, and so does
% anything else that has a name but no number.
%
% It places no anchor of its own. Under hyperref a sectioning command, starred
% or not, already places one at the title; a \label after it points there.
% \refstepcounter would place a second anchor here, below the title, and a
% link would lead below the heading. So the counter is stepped with
% the plain \stepcounter and the two things this needs are set here:
% \@currentlabel, which is what a \label records, and \@currentcounter,
% which is what zref-clever reads to give the label its reference type. The
% group keeps every change local.
%
% <text> is expanded to make \@currentlabel, and \label does not survive that
% expansion; so a \label inside <text> (for example a footnote's) is dropped
% from the label's text only. At its own place in the document it still makes
% its own label.
%
% A caller with no sectioning command before it -- for example a defined term
% in running text -- writes \phantomsection at its own place, so that a link
% has an anchor of its own to lead to.
%
% Redefine \flmLabelsWithDisplayText to change how such labels are made.
\ifcsname c@flm@internal@refcnt\endcsname\else\newcounter{flm@internal@refcnt}\fi
\ProvideDocumentCommand\flmLabelsWithDisplayText{m m}{%
  \begingroup
    \stepcounter{flm@internal@refcnt}%
    {\let\label\@gobble\protected@xdef\flm@lwdt@text{#1}}%
    \let\@currentlabel\flm@lwdt@text
    \def\@currentcounter{flm@internal@refcnt}%
    #2%
  \endgroup}
% flm.latex.def.zref-clever-type
% The zref-clever reference type of the counter \flmLabelsWithDisplayText
% steps, with every name prefix empty, so that a reference to such a label
% prints its text alone ("Introduction", not "Section Introduction").
% Every other label keeps the name its own type gives it. Redefine these
% settings, or the type of the counter, to put something before the text.
\zcsetup{countertype={flm@internal@refcnt=flm--internal--ref}}
\zcRefTypeSetup{flm--internal--ref}{Name-sg = , name-sg = , Name-pl = , name-pl = ,}
% flm.latex.def.labels
% \flmLabels{<\label commands>}: the labels of one place in the document.
% In text it runs them as written. On a line of a display-math
% environment, where amsmath keeps the last label of a line and loses the
% ones before it, it runs the first and records every further one as an alias:
% a line per alias is written to the .aux file after the last page and read on
% the next run, making the alias refer to what the first label refers to
% (for \ref and for \zcref alike). Redefine \flmLabels to change what a
% group of labels does; \def\flmLabels#1{#1} restores plain LaTeX behavior,
% amsmath's "Multiple \label's" message included.
\ProvideDocumentCommand\flmLabels{m}{%
  \ifmmode\expandafter\flm@labels@math\else\expandafter\@firstofone\fi{#1}}
% Reading the group: the first \label as written, the rest as aliases of it.
\def\flm@labels@math#1{\def\flm@lbl@first{}\flm@lbl@scan#1\flm@lbl@stop}
\def\flm@lbl@scan{\@ifnextchar\label\flm@lbl@take\flm@lbl@rest}
\def\flm@lbl@take\label#1{%
  \ifx\flm@lbl@first\@empty
    \def\flm@lbl@first{#1}%
    \label{#1}%
  \else
    \expandafter\flm@recordalias\expandafter{\flm@lbl@first}{#1}%
  \fi
  \flm@lbl@scan}
\def\flm@lbl@rest#1\flm@lbl@stop{#1}
% Collecting the aliases. A name is recorded once however often the line is
% typeset (amsmath measures a line before setting it).
\def\flm@aliaslist{}
\def\flm@recordalias#1#2{%
  \ifcsname flm@alias@#2\endcsname\else
    \expandafter\gdef\csname flm@alias@#2\endcsname{}%
    \g@addto@macro\flm@aliaslist{\flm@alias@entry{#2}{#1}}%
  \fi}
% Writing them, from the hook that runs after the last page has been shipped
% out: every \newlabel of that page has reached the .aux file by then, so the
% alias lines come after them. \AtEndDocument would run too early.
\def\flm@alias@write#1#2{\immediate\write\@auxout{\string\flm@aliaslabel{#1}{#2}}}
\AddToHook{enddocument/afterlastpage}{\let\flm@alias@entry\flm@alias@write\flm@aliaslist}
% Reading them back: the alias refers to what the original refers to. The
% assignments are global because the .aux file is read inside a group.
\def\flm@aliaslabel#1#2{%
  \ifcsname r@#2\endcsname
    \global\expandafter\let\csname r@#1\expandafter\endcsname\csname r@#2\endcsname
  \fi
  \ifcsname Z@R@#2\endcsname
    \global\expandafter\let\csname Z@R@#1\expandafter\endcsname\csname Z@R@#2\endcsname
  \fi}
% flm.latex.def.heading
% \flmHeading[<keys>]<sectioning command><star>{<title>}{<labels>}: a heading
% and the labels that refer to it, written on two lines for readability, as in
%   \flmHeading[toc=true]
%   \section*{Title}{\flmLabels{\label{ref1}}}
% It runs the sectioning command as written, then makes the labels refer to
% the heading: by its number when the class numbered it, by its title when it
% did not (a starred heading, or a level below secnumdepth). It finds out
% which without patching anything: a class numbers a heading through
% \refstepcounter, which records the counter's name in \@currentcounter.
%
% The optional argument is a key-value list. It is written here, on the hook,
% and never on the sectioning command: a starred sectioning command takes no
% optional argument, and giving one to \section* is not an error, it is
% silently wrong -- \section*[Short]{Long} takes one token as its heading,
% so the opening bracket alone becomes the heading and Short, the closing
% bracket and Long all run together into the following paragraph. Two keys
% are defined:
%
%   toc=true|false     whether the heading is listed in the table of contents.
%                      Absent, it is listed as LaTeX lists it. After a starred
%                      sectioning command with toc=true, which LaTeX never
%                      lists, \addcontentsline{toc}{<level>}{<entry>} adds it:
%                      <level> is the command's own name (section for
%                      \section), <entry> the alternative title if alt was
%                      written, else the title; under hyperref the starred
%                      command has already placed the anchor the line links
%                      to. On an unstarred heading toc=false is not honored:
%                      the heading stays listed and a warning says so.
%
%   alt={Short title}  an alternative title: what goes into the table of
%                      contents, the running head and the PDF bookmark, while
%                      the page keeps the full title. It is handed to the
%                      sectioning command as its own optional argument, so it
%                      does what \section[Short title]{Full title} does, and
%                      nothing else: a reference to an unnumbered heading still
%                      prints the full title. A starred heading does not hand
%                      it to the sectioning command; where it is written
%                      toc=true, alt is its line in the table of contents.
%
% What counts is whether the key was written, not what it holds. alt={} is an
% empty alternative title and is handed over as \section[]{Full title}, which
% is how LaTeX asks for an empty table-of-contents entry; no optional argument
% at all, and an empty one ([]), hand nothing over and leave the sectioning
% command exactly as it was. Writing alt with no value is LaTeX's own error,
% as is a key that is not defined here, and a value of toc other than true or
% false. The values and the record that alt was written are cleared at every
% call, so one heading's keys never reach the next. The recomposed LaTeX writes
% toc on every heading and alt where the heading was given an alternative
% title; both keys are also there for a document written by hand and for a
% user who redefines the hook.
%
% Redefine \flmHeading to change what a heading does. The block below uses
% expl3, LaTeX's programming layer, which is part of the format and needs no
% package; its names read only between \ExplSyntaxOn and \ExplSyntaxOff.
% \keys_define:nn declares the keys and \keys_set:nn reads a list of them;
% \bool_if:NTF asks whether the key was written. The three declarations are
% guarded on their own names, as every declaration in this file is, so that two
% of these style files can be loaded in one LaTeX run. The macros with @ in
% their names are internals, not a customization point: redefine \flmHeading,
% not its parts.
\ExplSyntaxOn
\ifcsname flm@heading@alt\endcsname\else\tl_new:N \flm@heading@alt\fi
\ifcsname flm@heading@given\endcsname\else\bool_new:N \flm@heading@given\fi
\ifcsname flm@heading@toc\endcsname\else\tl_new:N \flm@heading@toc\fi
\keys_define:nn {flm/heading}
  {alt .code:n = {\tl_set:Nn \flm@heading@alt {#1}
                  \bool_set_true:N \flm@heading@given} ,
   alt .value_required:n = true ,
   toc .choices:nn = {true, false} {\tl_set_eq:NN \flm@heading@toc \l_keys_choice_tl} ,
   toc .value_required:n = true}
\def\flm@heading@withalt#1#2#3{#2[#1]{#3}}
% The contents line of a starred heading: <level> then <entry>.
\def\flm@heading@addtoc#1#2{\addcontentsline{toc}{#1}{#2}}
% toc=false on an unstarred heading: the heading stays listed.
\def\flm@heading@notoc#1{\PackageWarning{flm}{toc=false~on~the~unstarred~
  \string#1~is~not~honored~by~the~default~\string\flmHeading;~the~heading~
  stays~in~the~table~of~contents}}
\ProvideDocumentCommand\flmHeading{O{} m s m m}{%
  \tl_clear:N \flm@heading@alt
  \bool_set_false:N \flm@heading@given
  \tl_clear:N \flm@heading@toc
  \keys_set:nn {flm/heading} {#1}%
  \def\@currentcounter{}%
  \IfBooleanTF{#3}
    {#2*{#4}%
     \str_if_eq:VnT \flm@heading@toc {true}
       {\edef\flm@heading@level{\cs_to_str:N #2}%
        \bool_if:NTF \flm@heading@given
          {\exp_args:NVV \flm@heading@addtoc \flm@heading@level \flm@heading@alt}
          {\exp_args:NV \flm@heading@addtoc \flm@heading@level {#4}}}}
    {\str_if_eq:VnT \flm@heading@toc {false} {\flm@heading@notoc{#2}}%
     \bool_if:NTF \flm@heading@given
       {\exp_args:NV \flm@heading@withalt \flm@heading@alt #2{#4}}
       {#2{#4}}}%
  \ifx\@currentcounter\@empty
    \flmLabelsWithDisplayText{#4}{#5}%
  \else
    #5%
  \fi}
\ExplSyntaxOff
% --- definitions written from the configuration of FLM features ---
% c@flmtheorem (the theorems feature's, written from its configuration): declared unless \c@flmtheorem is already defined, so a definition made earlier under that name is kept
\ifcsname c@flmtheorem\endcsname\else\newcounter{flmtheorem}\fi
% th@theoremlike (the theorems feature's, written from its configuration): declared unless \th@theoremlike is already defined, so a definition made earlier under that name is kept
\ifcsname th@theoremlike\endcsname\else\newtheoremstyle{theoremlike}{}{}{\itshape}{}{\bfseries}{.}{5pt plus 1pt minus 1pt}{}\fi
% theorem (the theorems feature's, written from its configuration): declared unless \theorem is already defined, so a definition made earlier under that name is kept
\ifcsname theorem\endcsname\else\begingroup\theoremstyle{theoremlike}\newtheorem{theorem}[flmtheorem]{Theorem}\endgroup\zcRefTypeSetup{theorem}{Name-sg={Theorem},name-sg={theorem},Name-pl={Theorems},name-pl={theorems},Name-sg-ab={Thm.},name-sg-ab={Thm.},Name-pl-ab={Thms.},name-pl-ab={Thms.}}\AddToHook{env/theorem/begin}{\zcsetup{countertype={flmtheorem=theorem}}}\fi
% proof (the theorems feature's, written from its configuration): declared unless \proof is already defined, so a definition made earlier under that name is kept
\ifcsname proof\endcsname\else\newenvironment{proof}[1][{Proof}]{\begin{proof}[{#1}]}{\end{proof}}\fi
% lemma (the theorems feature's, written from its configuration): declared unless \lemma is already defined, so a definition made earlier under that name is kept
\ifcsname lemma\endcsname\else\begingroup\theoremstyle{theoremlike}\newtheorem{lemma}[flmtheorem]{Lemma}\endgroup\zcRefTypeSetup{lemma}{Name-sg={Lemma},name-sg={lemma},Name-pl={Lemmas},name-pl={lemmas},Name-sg-ab={Lem.},name-sg-ab={Lem.},Name-pl-ab={Lems.},name-pl-ab={Lems.}}\AddToHook{env/lemma/begin}{\zcsetup{countertype={flmtheorem=lemma}}}\fi
% proposition (the theorems feature's, written from its configuration): declared unless \proposition is already defined, so a definition made earlier under that name is kept
\ifcsname proposition\endcsname\else\begingroup\theoremstyle{theoremlike}\newtheorem{proposition}[flmtheorem]{Proposition}\endgroup\zcRefTypeSetup{proposition}{Name-sg={Proposition},name-sg={proposition},Name-pl={Propositions},name-pl={propositions},Name-sg-ab={Prop.},name-sg-ab={Prop.},Name-pl-ab={Props.},name-pl-ab={Props.}}\AddToHook{env/proposition/begin}{\zcsetup{countertype={flmtheorem=proposition}}}\fi
% corollary (the theorems feature's, written from its configuration): declared unless \corollary is already defined, so a definition made earlier under that name is kept
\ifcsname corollary\endcsname\else\begingroup\theoremstyle{theoremlike}\newtheorem{corollary}[flmtheorem]{Corollary}\endgroup\zcRefTypeSetup{corollary}{Name-sg={Corollary},name-sg={corollary},Name-pl={Corollaries},name-pl={corollaries},Name-sg-ab={Cor.},name-sg-ab={Cor.},Name-pl-ab={Cors.},name-pl-ab={Cors.}}\AddToHook{env/corollary/begin}{\zcsetup{countertype={flmtheorem=corollary}}}\fi
% conjecture (the theorems feature's, written from its configuration): declared unless \conjecture is already defined, so a definition made earlier under that name is kept
\ifcsname conjecture\endcsname\else\begingroup\theoremstyle{theoremlike}\newtheorem{conjecture}[flmtheorem]{Conjecture}\endgroup\zcRefTypeSetup{conjecture}{Name-sg={Conjecture},name-sg={conjecture},Name-pl={Conjectures},name-pl={conjectures},Name-sg-ab={Conj.},name-sg-ab={Conj.},Name-pl-ab={Conjs.},name-pl-ab={Conjs.}}\AddToHook{env/conjecture/begin}{\zcsetup{countertype={flmtheorem=conjecture}}}\fi
% th@definitionlike (the theorems feature's, written from its configuration): declared unless \th@definitionlike is already defined, so a definition made earlier under that name is kept
\ifcsname th@definitionlike\endcsname\else\newtheoremstyle{definitionlike}{}{}{\itshape}{}{\bfseries}{.}{5pt plus 1pt minus 1pt}{}\fi
% definition (the theorems feature's, written from its configuration): declared unless \definition is already defined, so a definition made earlier under that name is kept
\ifcsname definition\endcsname\else\begingroup\theoremstyle{definitionlike}\newtheorem{definition}[flmtheorem]{Definition}\endgroup\zcRefTypeSetup{definition}{Name-sg={Definition},name-sg={definition},Name-pl={Definitions},name-pl={definitions},Name-sg-ab={Def.},name-sg-ab={Def.},Name-pl-ab={Defs.},name-pl-ab={Defs.}}\AddToHook{env/definition/begin}{\zcsetup{countertype={flmtheorem=definition}}}\fi
% remark (the theorems feature's, written from its configuration): declared unless \remark is already defined, so a definition made earlier under that name is kept
\ifcsname remark\endcsname\else\begingroup\theoremstyle{definitionlike}\newtheorem{remark}[flmtheorem]{Remark}\endgroup\zcRefTypeSetup{remark}{Name-sg={Remark},name-sg={remark},Name-pl={Remarks},name-pl={remarks},Name-sg-ab={Rem.},name-sg-ab={Rem.},Name-pl-ab={Rems.},name-pl-ab={Rems.}}\AddToHook{env/remark/begin}{\zcsetup{countertype={flmtheorem=remark}}}\fi
% idea (the theorems feature's, written from its configuration): declared unless \idea is already defined, so a definition made earlier under that name is kept
\ifcsname idea\endcsname\else\begingroup\theoremstyle{definitionlike}\newtheorem{idea}[flmtheorem]{Idea}\endgroup\zcRefTypeSetup{idea}{Name-sg={Idea},name-sg={idea},Name-pl={Ideas},name-pl={ideas},Name-sg-ab={Idea},name-sg-ab={Idea},Name-pl-ab={Ideas},name-pl-ab={Ideas}}\AddToHook{env/idea/begin}{\zcsetup{countertype={flmtheorem=idea}}}\fi
% question (the theorems feature's, written from its configuration): declared unless \question is already defined, so a definition made earlier under that name is kept
\ifcsname question\endcsname\else\begingroup\theoremstyle{definitionlike}\newtheorem{question}[flmtheorem]{Question}\endgroup\zcRefTypeSetup{question}{Name-sg={Question},name-sg={question},Name-pl={Questions},name-pl={questions},Name-sg-ab={Qtn.},name-sg-ab={Qtn.},Name-pl-ab={Qtns.},name-pl-ab={Qtns.}}\AddToHook{env/question/begin}{\zcsetup{countertype={flmtheorem=question}}}\fi
% claim (the theorems feature's, written from its configuration): declared unless \claim is already defined, so a definition made earlier under that name is kept
\ifcsname claim\endcsname\else\begingroup\theoremstyle{definitionlike}\newtheorem{claim}[flmtheorem]{Claim}\endgroup\zcRefTypeSetup{claim}{Name-sg={Claim},name-sg={claim},Name-pl={Claims},name-pl={claims},Name-sg-ab={Clm.},name-sg-ab={Clm.},Name-pl-ab={Clms.},name-pl-ab={Clms.}}\AddToHook{env/claim/begin}{\zcsetup{countertype={flmtheorem=claim}}}\fi
% observation (the theorems feature's, written from its configuration): declared unless \observation is already defined, so a definition made earlier under that name is kept
\ifcsname observation\endcsname\else\begingroup\theoremstyle{definitionlike}\newtheorem{observation}[flmtheorem]{Observation}\endgroup\zcRefTypeSetup{observation}{Name-sg={Observation},name-sg={observation},Name-pl={Observations},name-pl={observations},Name-sg-ab={Obs.},name-sg-ab={Obs.},Name-pl-ab={Obs.},name-pl-ab={Obs.}}\AddToHook{env/observation/begin}{\zcsetup{countertype={flmtheorem=observation}}}\fi
% problem (the theorems feature's, written from its configuration): declared unless \problem is already defined, so a definition made earlier under that name is kept
\ifcsname problem\endcsname\else\begingroup\theoremstyle{definitionlike}\newtheorem{problem}[flmtheorem]{Problem}\endgroup\zcRefTypeSetup{problem}{Name-sg={Problem},name-sg={problem},Name-pl={Problems},name-pl={problems},Name-sg-ab={Prob.},name-sg-ab={Prob.},Name-pl-ab={Probs.},name-pl-ab={Probs.}}\AddToHook{env/problem/begin}{\zcsetup{countertype={flmtheorem=problem}}}\fi
% exercise (the theorems feature's, written from its configuration): declared unless \exercise is already defined, so a definition made earlier under that name is kept
\ifcsname exercise\endcsname\else\begingroup\theoremstyle{definitionlike}\newtheorem{exercise}[flmtheorem]{Exercise}\endgroup\zcRefTypeSetup{exercise}{Name-sg={Exercise},name-sg={Exercise},Name-pl={Exercises},name-pl={Exercises}}\AddToHook{env/exercise/begin}{\zcsetup{countertype={flmtheorem=exercise}}}\fi
% --- user definitions ---
% flmUEkeptthm (the document's keptthm): defined unless the name is already taken, so a definition of it made before this file is loaded wins; \RenewDocumentEnvironment after loading this file changes it
\ProvideDocumentEnvironment{flmUEkeptthm}{+m +b}{\begin{theorem}[#1]\noproofref#2
\end{theorem}}{}
\endinput
